% !PRINTED 286
\weberpage{286}

By this mode of notation we are in a position to derive the group of the transformation equation at once from that of the division equation.

We first fix as domain of rationality the totality of the rational functions of $\varkappa^2$ with \gesperrt{rational numerical coefficients}.

By § 63, 3. the group of the division equation consists, in this domain of rationality, of all substitutions by which $\mu$, $\mu'$ are carried over into
\[
\begin{array}{c} \partial \mu - b \mu' \\ -c \mu + a \mu' \end{array}
\]
where $a$, $b$, $c$, $\partial$ are arbitrary integers taken according to the modulus $p$, whose determinant
\weq{9}{\Delta = a\partial - bc}
is not divisible by $p$; and the number of all these substitutions amounts to
\weq{10}{p(p-1)(p^2-1) \qquad [\S\ 63,\,(27)]}

But from this there results, by the mode of notation (4), (5), \gesperrt{the group of the transformation equation as consisting of all} the permutations which the symbol
\weq{11}{\left(z, \frac{\partial z - b}{-cz+a}\right) \ (\mathrm{mod}\ p)}
\gesperrt{expresses.}

We denote a substitution $(z, z')$, when
\weq{12}{z \equiv \frac{c+\partial z'}{a+bz'} \ (\mathrm{mod}\ p)}
holds, similarly as before by
\weq{13}{\begin{pmatrix} a, & b \\ c, & \partial \end{pmatrix},}
and obtain for the composition of two such substitutions, when
\[
z' \equiv \frac{c'+\partial' z''}{a'+b'z''} \ (\mathrm{mod}\ p)
\]
holds, the rule:
\weq{14}{(z,z')(z'z'') \equiv \begin{pmatrix} a, & b \\ c, & \partial \end{pmatrix}\begin{pmatrix} a', & b' \\ c', & \partial' \end{pmatrix} \equiv \begin{pmatrix} aa'+bc', & ab'+b\partial' \\ ca'+\partial c', & cb'+\partial\partial' \end{pmatrix} \ (\mathrm{mod}\ p),}
%CONT>>
